A aproximação de funções não lineares é um desafio recorrente em diversas áreas da engenharia, ciência e tecnologia, especialmente em problemas de modelagem, controle e previsão. Métodos baseados em lógica fuzzy têm se destacado por sua capacidade de lidar com incertezas e representar conhecimento de forma interpretável \cite{ross2010fuzzy, zadeh1965fuzzy}. Dentre esses métodos, o sistema fuzzy de Takagi-Sugeno (TS) \cite{takagi1985fuzzy} é amplamente utilizado devido à sua flexibilidade e precisão na modelagem de sistemas complexos.

O modelo Takagi-Sugeno permite a construção de regras fuzzy cujos consequentes podem ser funções matemáticas (constantes ou lineares), proporcionando maior capacidade de aproximação em relação aos sistemas Mamdani tradicionais \cite{jang1997neuro}. A escolha adequada das funções de pertinência, operadores fuzzy e métodos de ajuste dos consequentes é fundamental para o desempenho do sistema \cite{babuska1996fuzzy}.

Neste trabalho, é proposta a implementação de um sistema fuzzy Takagi-Sugeno, nas versões de ordem zero e primeira ordem, para a aproximação da função não linear $f(x) = e^{-x/5} \cdot \sin(3x) + 0.5 \cdot \sin(x)$ no intervalo $x \in [0, 10]$. São exploradas diferentes funções de pertinência (gaussianas, triangulares e trapezoidais) e técnicas de otimização dos consequentes, como gradiente descendente e Recursive Least Squares (RLS) \cite{haykin2002adaptive}. O desempenho do modelo é avaliado por meio do erro quadrático médio (RMSE) e da análise gráfica dos resultados obtidos.

Os experimentos realizados demonstram a eficácia do sistema fuzzy TS na aproximação de funções não lineares, destacando a importância da configuração dos parâmetros e da escolha dos métodos de ajuste para a obtenção de melhores resultados.